\begin{lemma}
\label{lemma-direct-sum-devissage}
Let $M$ be an $R$-module.  If $(M_{\alpha})_{\alpha \in S}$ is a direct sum
d\'evissage of $M$, then
$M \cong \bigoplus_{\alpha + 1 \in S} M_{\alpha + 1}/M_{\alpha}$.
\end{lemma}

\begin{proof}
By property (3) of a direct sum d\'evissage, there is an inclusion
$M_{\alpha + 1}/M_{\alpha} \to M$ for each $\alpha \in S$.  Consider the
map
$$
f : \bigoplus\nolimits_{\alpha + 1\in S} M_{\alpha + 1}/M_{\alpha} \to M
$$
given by the sum of these inclusions.
Further consider the restrictions
$$
f_{\beta} :
\bigoplus\nolimits_{\alpha + 1 \leq \beta} M_{\alpha + 1}/M_{\alpha}
\longrightarrow
M
$$
for $\beta\in S$. Transfinite induction on $S$ shows that the image of
$f_{\beta}$ is $M_{\beta}$. For $\beta=0$ this is true by $(0)$. If $\beta+1$
is a successor ordinal and it is true for $\beta$, then it is true for 
$\beta + 1$ by (3).  And if $\beta$ is a limit ordinal and it is true for
$\alpha < \beta$, then it is true for $\beta$ by (2). Hence $f$ is surjective
by (1).  

\medskip\noindent
Transfinite induction on $S$ also shows that the restrictions $f_{\beta}$
are injective. For $\beta = 0$ it is true. If $\beta+1$ is a
successor ordinal and $f_{\beta}$ is injective, then let $x$ be in the kernel
and write $x = (x_{\alpha + 1})_{\alpha + 1 \leq \beta + 1}$ in terms of its
components $x_{\alpha + 1} \in M_{\alpha + 1}/M_{\alpha}$.  By property (3) and
the fact that the image of $f_{\beta}$ is $M_{\beta}$ both
$(x_{\alpha + 1})_{\alpha + 1 \leq \beta}$ and $x_{\beta + 1}$ map to $0$.
Hence $x_{\beta+1} = 0$ and, by the assumption that the restriction
$f_{\beta}$ is injective also $x_{\alpha + 1} = 0$
for every $\alpha + 1 \leq \beta$.  So $x = 0$ and $f_{\beta+1}$ is injective.
If $\beta$ is a limit ordinal consider an element $x$ of the kernel. Then $x$
is already contained in the domain of $f_{\alpha}$ for some $\alpha < \beta$.
Thus $x = 0$ which finishes the induction.
We conclude that $f$ is injective since $f_{\beta}$ is for each $\beta \in S$.
\end{proof}
